
\begin{defn}\label{0013} Suatu
 \index{C*@aljabar-$C^*$}%
 \index{aljabar!$C^*$-}%
\df{aljabar-$C^\ast$} adalah aljabar Banach $A$ dengan involusi yang memenuhi
   \[ \norm{a^*a} = \norm{a}^2 \]
untuk setiap $a \in A$.  Sifat norma ini biasanya disebut
 \index{C*@syarat-$C^*$}%
\emph{syarat-$C^*$}. Norma aljabar yang memenuhi syarat ini adalah
 \index{C*@norma-$C^*$}%
 \index{norma!$C^*$-}%
\df{norma-$C^*$}.  Suatu
 \index{C*@subaljabar-$C^*$}%
 \index{subaljabar!$C^*$-}%
\df{subaljabar-$C^*$} dari aljabar-$C^*$ $A$ adalah subaljabar-$*\,$ tertutup dari~$A$.

Kita menyatakan kategori aljabar-$C^*$ dan homomorfisme-$*\,$ dengan $\cat{CSA}$.
 \index{kategori!$\cat{CSA}$ sebagai suatu}%
 \index{CSA@$\cat{CSA}$!kategori}%
Sepintas mungkin tampak aneh bahwa dalam kategori ini, yang objek-objeknya memiliki sifat aljabar sekaligus topologis,
morfismenya hanya disyaratkan mempertahankan struktur aljabar. Namun, ternyata
setiap morfisme dalam $\cat{CSA}$ secara otomatis kontinu---bahkan, kontraktif (lihat
proposisi~\futurexref{12.3.16}{00152171})---dan
 \index{morfisme bijektif!dapat dibalik!dalam $\cat{CSA}$}%
setiap morfisme injektif merupakan isometri (lihat proposisi~\futurexref{12.3.17}{00152181}).  Jelas bahwa
(seperti pada definisi~\ref{00109}) setiap morfisme bijektif dalam kategori ini merupakan isomorfisme. Dengan demikian, setiap homomorfisme-$*\,$ bijektif
merupakan isomorfisme-$*\,$ isometrik.
\end{defn}

\begin{exam} Ruang vektor $\C$ dari bilangan kompleks, dengan perkalian bilangan kompleks biasa
 \index{C*@aljabar-$C^*$!$\C$ sebagai suatu}%
 \index{C@$\C$!sebagai aljabar-$C^*$}%
dan konjugasi kompleks $z \mapsto \conj z$ sebagai involusi, merupakan aljabar-$C^*$
komutatif beridentitas.
\end{exam}

\begin{exam} Jika $X$ suatu ruang Hausdorff kompak, aljabar $\fml C(X)$ dari fungsi kontinu
 \index{C*@aljabar-$C^*$!$\fml C(X)$ sebagai suatu}%
 \index{C@$\fml C(X)$!sebagai aljabar-$C^*$}%
bernilai kompleks pada $X$ merupakan aljabar-$C^*$ komutatif beridentitas apabila involusinya
diambil sebagai konjugasi kompleks.
\end{exam}

\begin{exam} Jika $X$ suatu ruang Hausdorff kompak lokal, aljabar Banach
 \index{C*@aljabar-$C^*$!$\fml C_0(X)$ sebagai suatu}%
 \index{C@$\fml C_0(X)$!sebagai aljabar-$C^*$}%
$\fml C_0(X) = \fml C_0(X,\C)$ dari fungsi kontinu bernilai kompleks pada $X$ yang lenyap
di tak hingga merupakan aljabar-$C^*$ komutatif (tidak mesti beridentitas) apabila involusinya
diambil sebagai konjugasi kompleks.
\end{exam}

\begin{exam} Jika $(X,\mu)$ suatu ruang ukur, aljabar $\lfs\infty(X,\mu)$ dari
 \index{C*@aljabar-$C^*$!$L_\infty(S)$ sebagai suatu}%
 \index{L@$\lfs\infty(S)$!sebagai aljabar-$C^*$}%
fungsi terukur bernilai kompleks yang terbatas secara esensial pada $X$ (sekali lagi dengan konjugasi
kompleks sebagai involusi) merupakan aljabar-$C^*$.  (Secara teknis, tentu saja, anggota
$\lfs\infty(X,\mu)$ adalah kelas-kelas ekuivalensi fungsi yang berbeda pada himpunan-himpunan berukuran nol.)
\end{exam}

\begin{exam}\label{001305fa} Aljabar $\ofml B(H)$ dari operator linear terbatas pada ruang Hilbert
 \index{C*@aljabar-$C^*$!$\ofml B(H)$ sebagai suatu}%
 \index{B@$\ofml B(V)$!sebagai aljabar-$C^*$}%
$H$ merupakan aljabar-$C^*$ beridentitas.  Keluarga $\ofml K(H)$ dari operator kompak merupakan ideal-$*\,$ tertutup
dalam aljabar tersebut.
\end{exam}

\begin{exam}\label{001306} Sebagai kasus khusus dari contoh sebelumnya, perhatikan bahwa himpunan
 \index{C*@aljabar-$C^*$!M@$\mathbf M_n$ sebagai suatu}%
 \index{Ma@$\mathbf M_n = \M n\C$!sebagai aljabar-$C^*$}%
$\mathbf M_n$ dari matriks bilangan kompleks berukuran $n \times n$ merupakan aljabar-$C^*$ beridentitas.
\end{exam}

\begin{prop}\label{prop_K_is_clo_FR} Dalam ruang Hilbert $H$, ideal $\ofml K(H)$ dari operator kompak merupakan tutupan
ideal $\ofml{FR}(H)$ dari operator berperingkat hingga.
\end{prop}

Selama kira-kira satu dasawarsa sebelum Perang Dunia Kedua, sejumlah matematikawan Polandia terkemuka berkumpul secara teratur di kedai
kopi di Lwow untuk membahas hasil dan mengajukan masalah.  Mula-mula mereka bertemu di \emph{Caf\'e Roma}, kemudian di
\emph{Scottish Caf\'e}.  Pada 1935 Stefan Banach menyediakan buku catatan berjilid tempat mereka dapat menuliskan masalah-masalah tersebut.
Dari sinilah \emph{Scottish Book}~\cite{Mauldin:1981} yang kini terkenal bermula.  Pada 6 November 1936, Stanislaw Mazur
menuliskan Masalah 153, yang dikenal sebagai \emph{masalah aproksimasi}; pada intinya masalah itu menanyakan apakah proposisi sebelumnya
berlaku untuk ruang Banach.  Dapatkah setiap operator kompak pada ruang Banach diaproksimasi dalam norma oleh operator berperingkat hingga?
Sebagai hadiah bagi pemecah masalah ini, Mazur menawarkan seekor angsa hidup (tergolong salah satu hadiah paling murah hati yang dijanjikan dalam
\emph{Scottish Book}).  Masalah tersebut tetap terbuka selama beberapa dasawarsa, hingga pada 1972 Per Enflo memberikan jawaban negatif,
setelah berhasil menemukan ruang Banach refleksif separabel yang di dalamnya aproksimasi semacam itu gagal.  Dalam sebuah upacara yang diadakan
tahun itu di Warsawa, Enflo menerima angsa hidupnya dari Mazur sendiri.  Hasil tersebut diterbitkan dalam Acta
Mathematica~\cite{Enflo:1973} pada~1973.

\begin{exam} Operator diagonal $\diag(1,\frac12,\frac13,\dots)$ merupakan operator kompak pada ruang Hilbert~$l_2$.
\end{exam}











\section{Isometri Parsial}
\begin{defn}  Suatu elemen $v$ dari aljabar-$C^*$ $A$ disebut
 \index{parsial!isometri}%
\df{isometri parsial} jika $v^*v$ merupakan proyeksi dalam~$A$.  (Karena $v^*v$ selalu
swaadjoin, cukup disyaratkan bahwa $v^*v$ idempoten.)  Elemen $v^*v$ adalah
\df{proyeksi}
 \index{awal!proyeksi}%
 \index{proyeksi!awal}%
 \index{parsial!isometri!proyeksi awal dari suatu}%
\df{awal} (atau
 \index{tumpuan!proyeksi}%
 \index{proyeksi!tumpuan}%
 \index{parsial!isometri!proyeksi tumpuan dari suatu}%
\df{tumpuan}) dari~$v$, sedangkan $vv^*$ adalah \df{proyeksi}
 \index{akhir!proyeksi}%
 \index{proyeksi!akhir}%
 \index{parsial!isometri!proyeksi akhir dari suatu}%
\df{akhir} (atau
 \index{jangkauan!proyeksi}%
 \index{proyeksi!jangkauan}%
 \index{parsial!isometri!proyeksi jangkauan dari suatu}%
\df{jangkauan}) dari~$v$.  (Sebagai konsekuensi langsung dari proposisi berikutnya,
jika $v$ suatu isometri parsial, maka $vv^*$ memang merupakan proyeksi.)
\end{defn}

\begin{prop}\label{pi0111} Jika $v$ suatu isometri parsial dalam aljabar-$C^*$, maka
 \begin{enumerate}[label=\rm{(\alph*)}]
  \item $vv^*v = v$\,; dan
  \item $v^*$ merupakan isometri parsial.
 \end{enumerate}
\end{prop}

\begin{proof}[\emph{Petunjuk untuk bukti}] Misalkan $z = v - vv^*v$ dan tinjau $z^*z$.  \ns
\end{proof}

\begin{prop}\label{pi0114} Misalkan $v$ suatu isometri parsial dalam aljabar-$C^*$~$A$.  Maka proyeksi awalnya
$p$ adalah proyeksi terkecil (terhadap urutan parsial~$\preceq$
pada $\fml P(A)$\,) sedemikian sehingga $vp = v$, dan proyeksi akhirnya $q$ adalah proyeksi
terkecil sedemikian sehingga $qv = v$.
\end{prop}

\begin{prop}\label{pi0117} Jika $V$ suatu isometri parsial pada ruang Hilbert $H$ (yakni, jika $V$ suatu isometri parsial
dalam aljabar-$C^*$~$\ofml B(H)$\,), maka proyeksi awal $V^*V$ adalah
proyeksi dari $H$ pada $(\ker V)^\perp$, dan proyeksi akhir $VV^*$ adalah proyeksi
dari $H$ pada~$\ran V$.
\end{prop}

Berdasarkan hasil sebelumnya, $(\ker V)^\perp$ disebut
 \index{awal!ruang}%
 \index{ruang!awal}%
 \index{ruang!tumpuan}%
 \index{tumpuan!ruang}%
\df{ruang} \df{awal} (atau \df{tumpuan}) dari~$V$, sedangkan $\ran V$ kadang-kadang disebut
 \index{akhir!ruang}%
 \index{ruang!akhir}%
\df{ruang akhir} dari~$V$.

\begin{prop}\label{pi0121} Suatu operator pada ruang Hilbert merupakan isometri parsial jika dan hanya jika operator tersebut merupakan
isometri pada komplemen ortogonal dari kernelnya.
\end{prop}
